% Pass 1 cuts from Chapter 6 (Hybrid representations for similarity and label noise).
% Everything below was removed from the chapter file in pass 1: the old notes blocks (rewritten
% into takeaway + content-bullet form), 'Decision for you' bullets (now in DECISIONS.md),
% and agent-added LaTeX reverted to the 9763e83 text (tables, a lead sentence, edited \todo lines).


% ==== FROM: \chapter{Hybrid Representations for Similarity and Label Noise}\label{chap:similarity} ====
% ---- NOTES: Hybrid Representations for Similarity and Label Noise ----
% Apart from the preamble in Paragraph 1 below, if there is one, no text goes directly under the chapter heading. The chapter opens with its Introduction section, which carries the research questions, the status of the results and the bridge from Chapters \ref{chap:flaws} and \ref{chap:method}.
% Decision for you: does this chapter's paper get submitted before the thesis? If yes, the chapter gets a preamble on reused material (Paragraph 1 below). If no, there is no preamble and the Introduction's sentence on the status of the results covers it.
% Paragraph 1 (only if the answer above is yes): preamble on reused material
% - Same form as the Chapter \ref{chap:robustness} preamble: one short italic paragraph before the first section.
% - It names the manuscript with the same title and author order as its entry under "Research Papers" in the Summary of Research Contributions. That entry's working title and authors are not yet chosen; copy them once they are.
% - It says the figures, tables and text are identical or similar to that manuscript's.
% ---- END NOTES ----

% ==== FROM: \section{Introduction} ====
% ---- NOTES: Introduction ----
% Short: about one page. Three paragraphs, then the research questions as a list.
% Paragraph 1: what Chapters \ref{chap:flaws} and \ref{chap:method} leave open.
% - Chapter \ref{chap:flaws} showed where the similarity assumption fails on standard representations.
% - Example: ECFP4 bits compared by Tanimoto mix a few task-relevant features with many that are constant or irrelevant on the dataset.
% - Methods that trust a molecule's neighbours inherit that error. This includes finding noisy labels by neighbours (see "Noise detection by neighbours" in Chapter \ref{chap:flaws}; do not re-define it).
% - Chapter \ref{chap:method} built the adaptive hybrid representation: a budget of features picked from several base representations for one task and one model. It is an input, not a model.
% - This chapter asks whether that representation repairs what Chapter \ref{chap:flaws} found broken: similarity, noise detection, and learning with noisy labels.
% Paragraph 2: the original motivation (one or two sentences).
% - The hybrid started as a fix for similarity: select the task-relevant features first, then compute distances only on those.
% - Chapter \ref{chap:method} then used hybrids as ordinary model inputs. This chapter returns to the original motivation.
% Paragraph 3: the neighbour-based and label-noise methods under test (two or three sentences; name and cite them, do not re-define them)
% - They are defined in Chapter \ref{chap:background} and in "Noise detection by neighbours" in Chapter \ref{chap:flaws}. Point there.
% - Name three: nearest-neighbour editing \citep{Wilson1972}, filtering by classifiers trained on other folds \citep{Brodley1999}, and confident learning \citep{Northcutt2021}.
% - Each depends on the space it works in. The chapter tests whether a hybrid changes how well they work.
% List: research questions, one per results section, in the order of the results sections.
% - Distances and nearest-neighbour prediction: is predicting each molecule's activity from its five nearest neighbours more accurate in a hybrid than in the single representations it is built from?
% - Noise detection and estimation: are noisy labels easier to find, or their size easier to estimate, in a hybrid?
% - Learning with noisy labels: does a hybrid change how well noise mitigation (label corrections, cleaning) and noise-robust training methods work?
% Paragraph 4: status of the chapter (one sentence).
% - Planned as the second hybrid paper. Methods are final; results are preliminary.
% - TODO(results): one-sentence preview of the pattern found for each research question.
% ---- END NOTES ----

% ==== FROM: \section{Methods} ====
% ---- NOTES: Methods ----
% Paragraph 1: what this section covers (short signpost)
% - Three parts: similarity in a hybrid, the noise conditions used, and the three experiments on learning with noisy labels.
% - Every hybrid here comes from the search in Chapter 5. How a hybrid is built does not change.
% - Two things are new. One is a second score a candidate hybrid can be ranked on (Similarity in a hybrid). The other is what is done to the training labels.
%
% Paragraph 2: what the three experiments share
% - Datasets, scaffold-group-aware folds and metrics are as in Chapter 5 (Models and metrics). Which datasets each experiment uses goes in the table under Noise conditions.
% - Noise touches training labels only; Noise conditions states this. One clause here, no more.
% - Comparison: each downstream model's own top two single representations, ranked by that model's own score. Never a fixed list or one chosen fingerprint.
% - Reporting is as in Chapter 5 (Models and metrics). The recovery rate is defined once, in the Recovery rate subsubsection under Methods for learning with noisy labels.
%
% Paragraph 3: which labels the hybrid saw
% - In all three experiments the hybrid's features are rebuilt within each fold from that fold's training labels as the experiment presents them: noisy, clipped or uncleaned.
% - Which search settings (representations, budget, allocation method) the hybrid uses differs by experiment:
%   - Survey of twelve LNL methods: searched again in each fold and noise level on the noisy training labels, scored by the same neural network the methods use.
%     Budgets searched: 50, 100 and 250 features, with 2 to 3 inner folds. Chapter 5 searches eight budgets, from 10 to 2000 features.
%     If that search returns no score in a fold, the fold uses the neural network's settings from the clean-label search instead, and the results record which folds did. With no such settings stored, the fold is not scored.
%   - Labels clipped at a detection limit: taken from the search on the clean (unclipped) labels. Only the feature ranking and budget split inside each fold see the clipped labels.
%   - Cleaning real labels: taken from the search on the same real labels the experiment trains on. No clean labels exist for this case.
% - Why it matters: settings searched on clean labels hand the hybrid information a real user never has. The survey avoids this; the clipped-label experiment does not.
% - Decision for you: re-search the clipped-label experiment's hybrid settings on the clipped labels in each fold, as the survey does, or keep the clean-label settings and state it as a caveat. Re-searching makes the two experiments comparable and removes the clean-label information, at the cost of one hybrid search per clipping level and fold. Keeping it leaves the hybrid's numbers there resting on labels a user would not have.
%
% Paragraph 4: the three tasks kept apart
% - The chapter separates three tasks often treated as one: finding noisy labels, estimating how large the noise is, and improving the model's held-out score.
% - A method can do the first well and fail the third.
%
% Figure: one-panel flow. Left: dataset with injected noise, or none. Middle: the hybrid built per fold on the training labels, beside the top two single representations. Right: similarity score, and the held-out score after each LNL method.
% ---- END NOTES ----

% ==== FROM: \subsection{Similarity in a hybrid} ====
% ---- NOTES: Similarity in a hybrid ----
% This subsection is the one place the neighbour score and the weighted distance are defined. Other sections point here.
%
% Paragraph 1: two ways to use a hybrid for similarity
% - Ordinary hybrid: the search in Chapter \ref{chap:method} picks features that help a downstream model. Distances are then computed in that hybrid.
% - Similarity hybrid: the same search, but each candidate is ranked by how well each molecule's activity is predicted from its five nearest neighbours in the hybrid.
% - No downstream model is fitted while searching for a similarity hybrid.
% - Everything else is shared: importance methods, allocation methods, budgets and folds. Only the score that ranks candidates changes.
% - The neighbour score also sets each representation's share of the budget. So no internal evaluator is used.
% - For the budget split, the score is computed on a validation part carved from each fold's training molecules. The score that ranks whole candidates uses the fold's held-out molecules (Paragraph 3).
% - In the full runs, every representation in the dataset's set is a candidate. The number of representations (2, 3, 5 or all) is not searched for the similarity hybrid.
% - The search can rank each fold's representations by the neighbour score, on the same validation part of the training molecules. The full runs do not use this.
% - Decision for you: search the number of representations for the similarity hybrid too, or keep all of them. Searching it gives both hybrids the same search. Keeping all needs one sentence saying the similarity hybrid always draws on every representation.
% - Why it matters: Chapter \ref{chap:flaws} showed the representation decides which molecules count as neighbours. The similarity hybrid tests whether picking features for neighbour quality fixes that more directly than picking them for prediction.
%
% Paragraph 2: the distance inside a hybrid
% - Each feature is z-scored with the mean and standard deviation of the training molecules only (as in Chapter \ref{chap:method}; do not re-define).
% - The distance is Manhattan on the z-scored features. It is defined in "Distance metrics for continuous molecular descriptors" (Chapter \ref{chap:flaws}); point there.
% - Width is weight. Every z-scored feature adds one term of similar size. So a representation's pull on the distance is roughly its number of features divided by the budget.
% - Consequence: in a hybrid, the allocation method also weights the representations inside the distance.
% - The search's neighbour score uses the z-scored features with no importance weighting. Importance weighting is a separate comparison (next paragraph heading).
% - TODO(results): confirm Manhattan distance and one-over-distance neighbour weighting once the design experiment crossing distance, neighbour weighting and number of neighbours has run in full.
%
% Paragraph 3: the neighbour score
% - To rank whole candidates, the score is computed on the held-out molecules of each outer fold. Neighbours are other held-out molecules, not training molecules.
% - The held-out molecules' scaffold groups are split into up to 5 parts (fewer if there are fewer groups). Each molecule's 5 nearest neighbours come from the other parts, so from other scaffold groups.
% - Why: a neighbour with the same scaffold is an easy match and inflates the score.
% - Prediction, weighting each neighbour by one over its distance: $\hat{y}_i = \frac{\sum_{j\in N_5(i)} y_j / d_{ij}}{\sum_{j\in N_5(i)} 1/d_{ij}}$. A neighbour at distance 0 takes all the weight.
% - Regression score: Spearman $\rho(\hat{\mathbf{y}}, \mathbf{y})$.
% - Binary classification: $\hat{y}_i$ is the distance-weighted fraction of neighbours in class 1. The score is $2(\mathrm{AUROC} - 0.5)$.
% - Why the rescaling: chance is then 0 under both scores. This matters because the budget split divides by this score.
% - Related published measure: the modelability index, defined in Chapter \ref{chap:flaws}. It uses one nearest neighbour and a hard class match; this score keeps the graded neighbour fraction. Point there, one sentence.
% - The score measures how far the similar property principle holds in the hybrid. The principle is stated in "Molecular similarity and activity cliffs" (Chapter~\ref{chap:background}); point there in one clause, do not re-state or re-cite it.
%
% Paragraph 4: where the score is undefined
% - A fold is not scored, and is left out of the mean rather than counted as 0, when it has fewer than 10 held-out molecules, fewer than 2 scaffold groups, or constant labels or predictions.
% - Multiclass labels: AUROC is defined for 2 classes. This removes the PAM50 breast cancer subtypes (4 classes) and the trastuzumab HER2 multiclass dataset (3 classes).
% - The trastuzumab HER2 binary dataset: its folds are cut on edit-distance shells around the wild-type CDRH3. A held-out block can hold one shell, so no molecule has a neighbour from another group.
% - The similarity hybrid is not searched on these three datasets. It is searched on every other dataset in the thesis, antibody and other non-molecule datasets included.
%
% Figure: two pipelines sharing a first half (base representations -> importance -> allocation -> hybrid). Top: hybrid -> downstream model -> dataset metric. Bottom: hybrid -> 5 nearest neighbours from other scaffold groups -> distance-weighted mean activity -> Spearman or rescaled AUROC. Caption: same search, different score.
% ---- END NOTES ----
% (agent-edited \todo, reverted to the 9763e83 wording)
\todo[inline, color=yellow!40]{6.2.1: the search objective of predicting each molecule's activity from its five nearest neighbours. [DIAGRAM D6.1]}

% ==== FROM: \paragraph{Feature importance weighting.} ====
% ---- NOTES: Feature importance weighting ----
% Paragraph 1: what the weighted distance is, and where this chapter uses it (her paragraph below)
% - Her first sentence names noise detection as the purpose. The same weighted distance is also compared, weighted against unweighted, for predicting each molecule's activity from its nearest neighbours. Add that use.
% - Her formula ($w$ under Manhattan, $\sqrt{w}$ under Euclidean) is the general definition. Comparisons that compute the full distance matrix between molecules use exactly this: the matched power, rescaled to average 1.
% - Three experiments in this chapter multiply every feature by $\sqrt{w}$ whatever the distance. Whether that matches depends on the distance each one uses.
% - The noise-detection comparison finds neighbours under Manhattan, where the matching power is $w$ itself. So it does not match. It also takes the square root before rescaling to average 1, not after.
% - The comparison of weighted against unweighted features for predicting activity from nearest neighbours fits a nearest-neighbour model under Euclidean. That matches. Its part that builds features from each molecule's neighbours (their mean and spread of activity, and their distances) finds them under Manhattan. That does not match.
% - The survey of twelve LNL methods feeds the weighted hybrid features to the downstream model directly, with and without each LNL method. No distance is involved. The multiplier is $\sqrt{w}$, rescaled to average 1.
% - Decision for you: state $\sqrt{w}$ as the convention these experiments use, or move them to the matched power before the full runs. The first keeps the code and changes the definition sentence. The second keeps the sentence and changes one setting per experiment, so Manhattan then uses $w$.
% - The importance scores are computed on the training molecules of each fold only.
% - Weights are applied after z-scoring. Applied before it, z-scoring divides them back out and the distance is unchanged.
% - Why $w$ under Manhattan and $\sqrt{w}$ under Euclidean: either way each feature's term in the distance is multiplied by $w$ exactly once. For Minkowski with exponent $p$ the multiplier is $w^{1/p}$.
% - Negative importances set to 0 means a zero-importance feature drops out of the distance.
% - Average 1 keeps the weighted and unweighted distances on the same scale. Nearest-neighbour rankings are the same under any rescaling of $w$; only absolute cut-offs and kernel widths change.
%
% Paragraph 2: weighted against unweighted features for predicting each molecule's activity from its nearest neighbours
% - Question: does weighting each feature by its importance give a better distance between molecules than the raw features?
% - All base representations are joined into one matrix, not a hybrid. Importance is computed per fold on the training molecules, once per importance method.
% - Weight per feature: square root of its importance, normalised so the weights sum to 1 (not average 1 as above; nearest-neighbour rankings are the same under either). Features are standardised, then multiplied by their weights.
% - k-nearest-neighbour regression with 5, 10 and 20 neighbours, weighted against unweighted.
% - Second method: each held-out molecule's activity is the inverse-distance-weighted mean of its 10 nearest training molecules \citep{Dudani1976}, weighted against unweighted space.
% - Further: features selected by one importance method and weighted by another, at 25 features; and the 25 selected features with and without weights.
% - Comparator: the same 10-neighbour prediction on each single representation, unweighted.
% - Regression datasets only. Scored in $R^2$ and MAE, not in each dataset's own metric. The space is all representations joined, not a hybrid.
% - Decision for you: change the experiment before the full runs to score each dataset's own metric on the hybrid's features, or keep it and state both differences in the text. Changed, its numbers sit beside the other experiments' numbers. Kept, they compare only with each other.
% - Results are reported in "Results: distances and nearest-neighbour prediction".
% - It tests whether importance scores improve the distance this chapter uses. It supports the point in Chapter 7 that a smaller set of relevant features works better.
% - TODO(results): weighted minus unweighted, per dataset and number of neighbours.
% ---- END NOTES ----
% TODO: review, needs manual review

% ==== FROM: \subsection{Noise conditions} ====
% ---- NOTES: Noise conditions ----
% Paragraph 1: which conditions this chapter uses, and where they are defined.
% - The noise conditions are defined once in Chapter 2 (Section~\ref{sec:noise_injection}, Figure~\ref{fig:noise_conditions}). Point there; do not re-define them.
% - This chapter uses three of them for its three experiments on learning with noisy labels: censoring, Gaussian noise, and real labels with nothing injected.
% - Two more appear only in "Results: noise detection and estimation". Outlier noise and grouped-wider noise are used to grade detection. Grouped-shifted noise is used to grade estimation only, because every training label moves under it.
% - Table~\ref{tab:noise_conditions_lnl} gives, per condition, the levels, what it stands for, and which experiments and datasets use it. Point there.
% - Noise is added to the training labels of each outer fold only. Held-out labels are never changed.
% - The noise size, the censoring limit, and which scaffold families grouped-wider noise widens are all computed from that fold's clean training labels. Held-out labels never set any of them.
% - If the simulation of several sources per molecule with correlated errors is reported, state one exception where it is described. It sizes its noise from the standard deviation of all molecules' labels, held-out molecules included. Its held-out labels are still clean.
%
% Paragraph 2: censoring at a detection limit.
% - Defined in Section~\ref{sec:noise_injection}; the level is the fraction of training labels clipped. Point there.
% - The levels, and which of them each experiment runs, are in Table~\ref{tab:noise_conditions_lnl}. The 0 level is the negative control.
% - The clip has no random draw, so one draw per level is run.
% - The property the argument needs: the error is one-sided, a bias rather than a scatter, and its direction and its location (the limit) are known.
% - Stands for: assay range limits, e.g. "> 10 µM" recorded as 10 µM \citep{Svensson2025}. How common this is was given in Chapter 2 ("Noise in molecular data"); point there.
% - Regression only: a class label has no value to clip.
%
% Paragraph 3: Gaussian noise.
% - Defined in Section~\ref{sec:noise_injection}. Point there.
% - Levels are in Table~\ref{tab:noise_conditions_lnl}. Three noise draws per level.
% - Stands for: the random error of repeating a measurement. It is also the standard assumption in noise-robustness studies (Chapter 2, "Noise in molecular data").
% - Used by the survey of twelve LNL methods.
% - Why it is the hardest case to correct is argued once, in "By noise type". Point there.
% - On a classification dataset the level is instead the fraction of training labels flipped to the wrong class (Table~\ref{tab:noise_conditions_lnl}).
%
% Paragraph 4: outlier, grouped-wider and grouped-shifted noise (used for detection and estimation only).
% - All three are defined, with their parameters, in Chapter 2 (Section~\ref{sec:noise_injection}). Point there; one line each here on what they stand for.
% - What differs from Chapter 2: outlier noise is run with 1\%, 5\% and 10\% of training molecules affected, not only the 10\% of Chapter 2.
% - Outlier stands for annotation and transcription errors \citep{Kalliokoski2013, Hampel2001}.
% - Grouped-wider and grouped-shifted stand for differences between laboratories or sources \citep{Bentz2013}.
% - Both grouped conditions need a scaffold per molecule, so they run only on datasets with SMILES.
% - None of the three has an LNL result in this chapter. "By noise type" says so; do not repeat it here.
%
% Paragraph 5: real labels, nothing injected.
% - Labels are used as they come. No list of wrong labels exists.
% - Success is therefore only a better held-out score after cleaning the training labels.
% - A win shows cleaning helped the model, not that it found errors. It may have removed hard molecules instead.
% - QM9 is the control: its labels come from a quantum-chemical calculation, with no experimental error \citep{Ramakrishnan2014}. A cleanup that helps on QM9 is not removing measurement error.
% ---- END NOTES ----
% TODO: review, needs manual review (table and caption)
\begin{table}[htbp]
\centering
\caption{The noise conditions used in this chapter, one row per condition. Each condition is
defined in Section~\ref{sec:noise_injection}. A level is a multiple of the standard deviation of
the fold's clean training labels, except for censoring (fraction of training labels clipped) and
for classification datasets (fraction of training labels flipped to the wrong class). Noise is
added to training labels only.}
\label{tab:noise_conditions_lnl}
\small
\setlength{\tabcolsep}{4pt}
\begin{tabular}{p{2.0cm}p{2.9cm}p{3.8cm}p{2.7cm}p{1.8cm}}
\toprule
\textbf{Condition} & \textbf{Stands for} & \textbf{Levels} & \textbf{Used by} & \textbf{Datasets} \\
\midrule
Censoring & Assay range limits \citep{Svensson2025} & 0, 0.10, 0.20, 0.25, 0.30, 0.40 of training labels clipped; the survey and the detection results run 0, 0.10, 0.25, 0.40. Regression only & Correcting clipped labels; survey of twelve LNL methods; noise detection &
% TODO(results): datasets each experiment covers in the full runs
\\
Gaussian & Random error of repeating a measurement & 0, 0.2, 0.3, 0.5, 0.75, 1.0, 1.5; classification: 0, 0.05, 0.10, 0.20, 0.30, 0.40 flipped & Survey of twelve LNL methods &
% TODO(results): datasets each experiment covers in the full runs
\\
Real labels, nothing injected & Every source, as recorded & None & Cleaning real labels &
% TODO(results): datasets each experiment covers in the full runs
\\
Outlier & Annotation and transcription errors \citep{Kalliokoski2013, Hampel2001} & As Gaussian, with 1\%, 5\% or 10\% of training molecules affected & Noise detection &
% TODO(results): datasets each experiment covers in the full runs
\\
Grouped-wider & Differences between laboratories \citep{Bentz2013} & As Gaussian; datasets with SMILES only & Noise detection &
% TODO(results): datasets each experiment covers in the full runs
\\
Grouped-shifted & Differences between laboratories \citep{Bentz2013} & As Gaussian; datasets with SMILES only & Noise estimation &
% TODO(results): datasets each experiment covers in the full runs
\\
\bottomrule
\end{tabular}
\end{table}

% ==== FROM: \subsection{Methods for learning with noisy labels} ====
% ---- NOTES: Methods for learning with noisy labels ----
% Paragraph 1: what this subsection covers, and the two kinds of LNL method
% - Learning with noisy labels (LNL), noise mitigation and noise-robust training are defined in Section~\ref{sec:lnl_background}. Point there in one clause; do not define them again here.
% - The one point that is new here: noise-robust training never has to say which direction a label is wrong in. Noise mitigation does, per molecule.
% - Three experiments, one per question: a known, one-sided defect (labels clipped at a detection limit); a survey of twelve LNL methods under Gaussian noise; cleaning real labels with nothing injected.
% - The first and third test noise mitigation. The survey covers both kinds.
% - Feature spaces, and which labels each hybrid was built on, are set out in Methods, earlier in this chapter. Say only what each experiment adds.
%
% Paragraph 2: the existing paragraph below this notes block
% - It tells the route of the exploratory study (Section~\ref{sec:lnl_motivation}): removal first, then weighting and smoothing.
% - The three experiments here do not follow that route. Removal is back: four of the twelve surveyed methods remove molecules, and removal is one of seven cleanups.
% - Decision for you: cut that paragraph to one bridge sentence from the exploratory study to the three experiments, or move it whole to Section~\ref{sec:lnl_motivation}. A bridge sentence keeps a link here but drops the removal-to-smoothing story; moving it keeps the story where the study is told, and this subsection then opens on Paragraph 1.
%
% Paragraph 3: how each experiment is scored
% - Clipped labels: recovery rate, defined in the Recovery rate subsubsection below. There is a clean score to recover towards. Do not restate the formula here.
% - For a model comparator, the denominator is the clean-to-clipped gap of the model it replaces.
% - Survey: each method's held-out score against the plain network on the same input and noise level.
% - Real labels: no clean score exists. The headline is folds where the cleaned model beats the uncleaned one, out of folds run. Mean difference and spread across folds go beside it.
%
% Paragraph 4: correcting labels clipped at a detection limit (setup)
% - Training labels above a limit $c$ are recorded as $c$. How the true limit and the clipping levels are set, and why only regression datasets are used, is in Noise conditions. Point there; do not repeat it.
% - The treatments never see the true limit. They read it off the recorded training labels: the largest label, when at least 1\% of training molecules and at least 2 molecules sit on it.
% - Clipped molecules are those whose recorded label equals the limit. This is what a user without the true values can see.
% - Reference prediction $p_i$: out-of-fold boosted-tree prediction on the fold's training molecules, with inner folds keeping scaffold groups together. It is computed in the same feature space the treated model is fitted in.
% - Residual spread $s$: standard deviation of $y_i - p_i$ over molecules not at the limit. Including clipped molecules would mix in the distance the clip moved them.
% - Basis of the corrections, the mean of a normal above a limit: $\mathbb{E}[Y \mid Y > c] = \mu + \sigma\frac{\phi(a)}{1-\Phi(a)}$, $a = (c-\mu)/\sigma$. The censored (Tobit) model is built on this \citep{Tobin1958}.
%
% Paragraph 5: the label treatments and the model comparators
% - Clean labels (upper reference) and clipped labels as recorded (lower reference).
% - Per-molecule lift: each clipped label set to $c + s\frac{\phi(a_i)}{1-\Phi(a_i)}$, $a_i = (c-p_i)/s$. It starts from $c$, not $p_i$, so it exceeds the truncated-normal mean by $c - p_i$.
% - Shuffled control: the same per-molecule lifts dealt to other clipped molecules, three random orders, each reported on its own. It separates "labels were raised" from "each molecule got its own amount".
% - Constant lift: every clipped label raised by one amount, the mean of the per-molecule lifts above $c$.
% - Truncated-normal mean per molecule: $p_i + s\frac{\phi(a_i)}{1-\Phi(a_i)}$, never below $c$.
% - Drop the clipped molecules.
% - Winsorising: every training label moved to $p_i + \mathrm{clip}(y_i - p_i, -k\hat{s}, k\hat{s})$, with $\hat{s} = 1.4826 \times$ the median $|y_i - p_i|$ and $k \in \{1, 1.25, 1.5\}$. Labels within $k\hat{s}$ of their prediction do not move.
% - Model changes are comparators, not proposals: the thesis changes the input, not the model. Each is paired with the model it replaces, on the same clipped labels and features. The censored-likelihood model is also paired against its plain version on the constant lift.
%   - gradient boosting with a censored likelihood that treats clipped labels as lower bounds;
%   - a random forest with absolute-error splits;
%   - gradient boosting with the Huber loss.
% - Feature spaces: those set out in Methods; nothing added.
% - TODO(results): recovery rate per treatment, model, feature space and clipping level; and whether dropping clipped molecules does better or worse than keeping them.
%
% Paragraph 6: a survey of twelve LNL methods
% - Question: does changing the input from a single representation to the hybrid change which methods work, and how well?
% - Gaussian noise on training labels at the levels in the table under Noise conditions, with level 0 as the clean reference.
% - One plain fully connected network is the base for every method: three hidden layers of 256, 128 and 64 units, 200 epochs, dropout 0.2.
% - Inputs: the top two single representations, the hybrid, and the importance-weighted hybrid (Similarity in a hybrid). Each input is run with the plain network and with each method.
% - Two methods (heteroscedastic GP, evidential regression) exist for continuous labels only; on a classification dataset ten methods run.
% - That most of these methods were built for image classification with flipped labels is said once, in Section~\ref{sec:lnl_motivation}. One clause pointing back there; do not restate it. Here say only that they are adapted to continuous labels. The reasons they may not transfer are in Chapter~\ref{chap:discussion}.
% - TODO(results): per noise level, how many methods beat the plain network on each input; and whether the input moves the score more than the method.
%
% Paragraph 7: cleaning real labels, nothing injected
% - Per fold: fit on the uncleaned labels and score; clean the training labels only; refit the same model and score again.
% - What a win does and does not show (it may have removed hard molecules, not wrong labels) is stated once in Noise conditions. Point there.
% - Suspicion scores, two: neighbour residual (distance of a label from the mean label of its 10 nearest neighbours, neighbours from other scaffold groups only; formula in Noise magnitude estimation) and cross-validation residual (distance from an out-of-fold prediction).
% - Seven cleanups, one line each in the text:
%   - remove the flagged molecules;
%   - down-weight molecules by suspicion;
%   - move every label toward its neighbours' mean, in proportion to suspicion;
%   - move only flagged labels toward their neighbours' mean;
%   - move labels toward an out-of-fold model prediction;
%   - move labels toward a Gaussian-process mean;
%   - clip labels to the range carried by similar training molecules.
% - The last three need continuous labels and are not run on classification datasets.
% - Remove and the flagged-only neighbour move act on a share of molecules flagged: 1\%, 5\% or 20\%. On classification datasets the proportional neighbour move also acts on the flagged share.
% - The others act on every molecule in proportion to its suspicion, so a strength is swept instead. Down-weighting: most-to-least suspect weight ratio 2:1, 4:1 or 10:1. The three moves: a quarter, half or all of the way to the target.
% - The range clip reads no suspicion score. The range is the 10th to 90th percentile of labels among training molecules that share its leaves in a quantile regression forest. It runs once.
% - Shuffled-score control: the same cleanups with suspicion scores dealt at random, three orders, each reported on its own. It separates "which molecules" from "how many".
% - Feature spaces: those set out in Methods, plus the hybrid ranked on the similarity score (Similarity in a hybrid). Each is built once per fold on the uncleaned labels.
% - QM9 is reported separately as the control; Noise conditions says why. One clause here.
% - TODO(results): folds won per cleanup, suspicion score and feature space.
%
% Paragraph 8: how the four subsubsections below relate to these experiments
% - Of their content, the three experiments use only: the neighbour-residual score, the cross-validation residual, proportional smoothing, inverse weighting, and the recovery rate.
% - The rest (neighbourhood-spread filtering, LOF, the list of removal methods, the probabilistic estimators, threshold and confidence smoothing, exponential and soft weights) describes the exploratory study (Section~\ref{sec:lnl_motivation}).
% - Decision for you: keep all four subsubsections here and trim them to what the three experiments use, or move the unused material to Section~\ref{sec:lnl_motivation}. Keeping them here holds all estimator definitions in one place, but most of it describes a study reported in another chapter. Moving it shortens this chapter; the Recovery rate subsubsection stays here either way, because all three experiments use it.
%
% Table~\ref{tab:lnl_methods}, after the existing paragraph below: the twelve LNL methods of the survey in three groups, one plain line each. Cite it from Paragraph 6.
% Figure: true label against recorded label for the clipped-label experiment, with the pile-up at $c$ and arrows for the lift treatments.
% ---- END NOTES ----
% TODO: review, needs manual review (whole table: grouping, one-line descriptions and caption)
\begin{table}[htbp]
\centering
\caption{The twelve LNL methods in the survey, in three groups. One row is one method; every method uses the same plain fully connected network as its base. Small-loss selection is counted as noise mitigation, not noise-robust training, because this implementation uses its loss ranking only once, to delete the flagged molecules before a fresh network is fitted. Per-epoch down-weighting is counted as noise-robust training because its weights change during training. The two methods that model the noise run on continuous labels only.}
\label{tab:lnl_methods}
\small
\begin{tabular}{p{0.3\textwidth}p{0.62\textwidth}}
\toprule
\textbf{Method} & \textbf{What it does} \\
\midrule
\multicolumn{2}{l}{\textit{Noise mitigation: remove suspected molecules, then refit}} \\
Small-loss selection \citep{Han2018} & Trains briefly on all molecules, keeps those with the smallest training loss, and fits a new network on them. \\
Confident learning \citep{Northcutt2021} & Removes molecules on which a small ensemble of networks has a large error and agrees with itself, then refits. \\
Removal by distance to neighbours & Removes molecules whose label is far from the mean label of their nearest neighbours in the feature space, then refits. \\
Staged removal \citep{zhao2025} & Fits, removes the molecules with the largest training loss, and repeats over several stages before a final fit. \\
\midrule
\multicolumn{2}{l}{\textit{Noise-robust training: change the training loop}} \\
Per-epoch down-weighting & After each epoch, lowers the weight of molecules whose loss is above the median; the weights carry over to the next epoch. \\
Curriculum by confidence \citep{Bengio2009} & After a warm-up, trains first on the molecules with the smallest loss and adds the rest in stages. \\
DivideMix \citep{li2020dividemix} & Two networks; each splits the molecules into likely clean and likely noisy by fitting a two-component Gaussian mixture to its losses, and each network trains on the split made by the other. \\
Meta-Weight-Net \citep{shu2019metaweightnet} & A small second network maps each molecule's loss to a training weight and is updated to lower the loss on held-out validation molecules. \\
MentorNet \citep{jiang2018mentornet} & A second network sets each molecule's weight from its current loss, its recent losses and its distance to its neighbours. \\
Contrast to Divide \citep{zheltonozhskii2021contrast} & Perturbs each molecule's features several times; molecules with a high loss and predictions that barely move are down-weighted. \\
\midrule
\multicolumn{2}{l}{\textit{Noise-robust training: model the noise directly}} \\
Heteroscedastic GP \citep{liu2020} & A Gaussian process predicts the label; a small network predicts each molecule's noise variance. \\
Deep evidential regression \citep{amini2020deep} & Predicts a distribution over the mean and variance of each label; the loss penalises confident wrong predictions. \\
\bottomrule
\end{tabular}
\end{table}
% TODO(lit): per-epoch down-weighting has no verified source. The one candidate, \citet{zhang2025}, analyses meta-learning reweighting and proposes a different method, not a per-epoch decay of high-loss weights. Find the source, or keep the row without a citation.

% ==== FROM: \subsubsection{Sample removal methods} ====
% ---- NOTES: Sample removal methods ----
% Paragraph 1: what this study was and where removal sits
% - Every method here flags training molecules and deletes them before refitting. That makes them noise-mitigation methods, a subset of learning with noisy labels (LNL).
% - This was the exploratory study behind the move to weighting and smoothing. The opening paragraph of "Methods for learning with noisy labels" states that move; do not repeat it here.
% - Data: the QM9 HOMO-LUMO gap, 10,000 molecules, one scaffold split. Training and validation molecules were pooled for training.
% - Noise: Gaussian noise added to every training label, standard deviation 0.3 and 0.6 in the label's units (eV). Test labels stay clean.
% - Feature spaces: ECFP4, PDV, MHG-GNN, and hybrids built from them.
% - Cite only the single-representation rows. Some hybrid rows were built with an allocation method that is no longer part of the thesis.
% - Model after removal: a random forest with 100 trees. Score: $R^2$ on the test molecules, and the recovery rate defined in "Recovery rate".
% - Each neighbour-based method was run under eight distances (Euclidean, Manhattan, cosine, Chebyshev, Canberra, correlation, Bray-Curtis, Mahalanobis), with and without importance weighting of the features.
% - Two settings are undefined and were dropped. Bray-Curtis is dropped at both weightings: it needs non-negative features, and the features were standardised first. Mahalanobis is run unweighted only: it re-estimates the covariance, so per-feature weights cancel.
% - Molecule count: 10,000 QM9 molecules, the same study as Section~\ref{sec:lnl_motivation} (the confirmation report gives 10,000 for this study). The 5,000-molecule runs in the notes to "Noise detection by neighbours" in Chapter~\ref{chap:flaws} are a separate weighting and smoothing study. Do not merge the two setups.
% - Whether this subsubsection stays in this chapter is asked once, in the notes to "Methods for learning with noisy labels" above (Paragraph 8).
%
% Paragraph 2: k-nearest-neighbour filtering (the existing prose matches the code)
% - Standard deviation of the labels of each molecule's $k$ nearest neighbours; flag above the 90th percentile; $k \in \{5, 10, 20\}$.
% - Consequence worth one sentence: a 90th-percentile cut flags 10% of training molecules every time, whatever the noise level.
% - Your first sentence ("Distance-based approaches rely on the assumption that similar molecules have similar properties") re-states the similar property principle from "Molecular similarity and activity cliffs" (Chapter~\ref{chap:background}), with a different citation. Point back to that section, or cite \citet{Johnson1990} as Chapter~\ref{chap:background} does.
% - The sentence beginning "Presumably" can state the idea directly and cite its origin. Editing a training set by its nearest neighbours goes back to \citep{Wilson1972}.
%
% Paragraph 3: Local Outlier Factor (the existing prose is out of date)
% - What the code does: LOF with 20 neighbours, flagging the 10% of training molecules with the highest LOF score. There is no "LOF greater than 2" rule.
% - LOF sees only the features, never the labels. It flags molecules that sit in sparse regions of the representation, not molecules with unusual labels.
% - The same is true of the Mahalanobis detector in the table below. Say this once, for both.
%
% Paragraph 4: the other removal methods (the one-sentence list becomes the table)
% - One lead-in sentence, then the table. Rename "AC detection" to "activity-cliff detection".
% - Co-teaching \citep{Han2018} and DivideMix \citep{li2020} are noise-robust training methods in their original form. Both train deep networks on a selected subset instead of deleting data.
% - Here each was cut down to its selection step, run with random forests, and its rejected molecules deleted. State this so a reader does not expect the original methods.
% - Cross-validation filtering follows \citep{Brodley1999}. In the current code the inner folds are consecutive blocks of the training set, not scaffold groups (unchecked whether the original runs split the same way).
%
% Paragraph 5: why removal does not suit Gaussian noise (one or two sentences, then point on)
% - Under Gaussian noise every training label is wrong by some amount. Any flagged molecule is a "true" detection, so detection precision is 1 by construction and cannot rank methods.
% - The useful question becomes how wrong each label is, which "Noise magnitude estimation" takes up.
% - Decision for you: bring back your own first-person paragraph on why removal was dropped for Gaussian noise ("Initially, I attempted to adapt noise mitigation techniques..."), deleted before this draft. Restored here, the pivot paragraph under "Methods for learning with noisy labels" repeats it and should be cut to one bridge sentence. Left out, that pivot paragraph is the only statement of the move.
% - Decision for you: report the exploratory removal numbers, or keep this subsubsection as motivation only. Reported, it needs the recovery rate per method at noise SD 0.3 and 0.6 eV, and the table below gains two columns. Motivation only, the table stays as it is.
% - TODO(results): if reported, the recovery rate of the best and the median removal method at noise SD 0.3 and 0.6 eV, single representations only.
%
% ---- END NOTES ----
% TODO: review, needs manual review
Table~\ref{tab:removal_methods} lists every removal method in the exploratory study, what it scores for each training molecule, and which molecules it flags.

% TODO: review, needs manual review
\begin{table}[htbp]
\centering
\small
\caption{Sample removal methods in the exploratory study. One row per method. ``Top 10\%'' means the 10\% of training molecules with the highest score. SD is the standard deviation of the training labels.}
\label{tab:removal_methods}
\begin{tabular}{p{2.6cm}p{5.6cm}p{3.2cm}p{1.1cm}}
\toprule
Method & Score per training molecule & Flagged & Uses labels \\
\midrule
$k$-nearest-neighbour filtering \citep{Wilson1972} & Standard deviation of the labels of the $k$ nearest neighbours, $k \in \{5, 10, 20\}$ & Top 10\% & Yes \\
\addlinespace
Local Outlier Factor \citep{Breunig2000} & Local density against that of its 20 nearest neighbours & Top 10\% & No \\
\addlinespace
Distance-weighted prediction & Error of the inverse-distance-weighted mean of the 15 nearest neighbours' labels & Error above 2 SD & Yes \\
\addlinespace
Neighbourhood consensus & Difference from the mean label of the 10 nearest neighbours & Difference above 2 SD and neighbours' spread below 1 SD & Yes \\
\addlinespace
Activity-cliff detection & Membership of a pair in the closest 10\% of pairs by distance and the largest 10\% by label difference & Every molecule in such a pair & Yes \\
\addlinespace
Mahalanobis & Distance from the centre of the standardised features & Top 10\% & No \\
\addlinespace
Cross-validation disagreement \citep{Brodley1999} & Out-of-fold error of a random forest (50 trees), 5 inner folds & Top 10\% & Yes \\
\addlinespace
Bootstrap ensemble & Error of the mean out-of-bag prediction of 10 bootstrap random forests (50 trees each) & Top 10\% & Yes \\
\addlinespace
Co-teaching, selection step only \citep{Han2018} & Two random forests each pick the other's 70\% lowest-error molecules, for 20 rounds & Rejected at least 12 of 40 times & Yes \\
\addlinespace
DivideMix, selection step only \citep{li2020} & Two-component Gaussian mixture on the out-of-bag errors of one random forest (100 trees) & Higher-error component & Yes \\
\addlinespace
Label outlier baseline & Distance of the label from the label mean & Above 3 SD & Yes \\
\addlinespace
Random forest error baseline & Out-of-bag error of a random forest & Top 10\% & Yes \\
\bottomrule
\end{tabular}
\end{table}

% ==== FROM: \subsubsection{Noise magnitude estimation} ====
% ---- NOTES: Noise magnitude estimation ----
% Paragraph 1: why score how noisy each label is, not only whether it is noisy
% - A score per molecule lets later steps down-weight or smooth a label in proportion. Removal is all or nothing.
% - Filtering suspect labels was built for classification. \citet{Martin2021} adapt fourteen classification noise filters to continuous labels and find they help regression models. Cite it for the general idea only.
% - Wording in the first sentence: estimating noise is not itself noise mitigation. It feeds noise mitigation (smoothing, weighting, removal). Say "the step that noise mitigation builds on".
% Paragraph 2: scores from each molecule's nearest neighbours in the representation space
% - State the number of neighbours: 10 per molecule. Where a molecule has repeat measurements, none of its repeats count as its neighbours.
% - Neighbour residual: distance of the label from the mean label of its 10 nearest neighbours, divided by the label standard deviation of the whole dataset. It is the continuous-label form of editing a training set by neighbour disagreement \citep{Wilson1972}.
% - Local variance ratio: the same distance, divided by the standard deviation of the neighbours' labels.
% - The prose writes this denominator as std + epsilon. What is computed is a floor: the neighbours' standard deviation, but never below 1% of the dataset's label standard deviation.
% - Activity-cliff-aware score: spell out "activity cliff" instead of "AC". An outlying label among neighbours that agree suggests noise. Among neighbours that disagree it suggests an activity cliff. The formula in the prose matches. Cite \citep{Maggiora2006, Stumpfe2012, Dablander2023}.
% - Distance-weighted neighbour residual, missing from the prose: the neighbour residual with each neighbour weighted by one over its distance.
% Paragraph 3: scores from models that never saw the molecule
% - Cross-validation residual: absolute error of an out-of-fold random forest prediction (100 trees), 5 folds, scaffold groups kept together, divided by the label standard deviation. Cite \citep{Brodley1999} for flagging labels that models trained on other data cannot reproduce.
% - Ensemble disagreement: 10 random forests of 50 trees, each with a different seed, predicting out of fold. Score = spread of their predictions + absolute error of their mean, divided by the label standard deviation.
% - Cite \citep{Brodley1999} for consensus filtering and \citep{Dietterich2000} for ensembles.
% - Out-of-bag residual, missing from the prose: absolute error of a 200-tree random forest's out-of-bag prediction. With repeat measurements, whole groups of repeats are held out together instead.
% - Combined score, missing from the prose: the neighbour residual, local variance ratio and cross-validation residual, each rescaled to 0-1, then averaged.
% Paragraph 4: scores from models that predict their own uncertainty
% - All of these are fitted out of fold, 5 folds, so no molecule is scored by a model that trained on it.
% - Gaussian process: the prose is out of date. One learned noise level is shared by every molecule, so it cannot rank molecules. The score is the out-of-fold predictive standard deviation plus the out-of-fold absolute residual \citep{gauche}.
% - Heteroscedastic Gaussian process, missing from the prose: a Gaussian process for the mean and a small network predicting each molecule's noise level. The score is that per-molecule noise level alone.
% - Heteroscedastic network: outputs a mean and a variance for each molecule \citep{Nix1994}. Trained with the beta-NLL loss at beta = 0.5 \citep{Seitzer2022}; the formula in the prose matches.
% - Conformal prediction: split conformal prediction for regression \citep{Lei2018}; tutorial \citep{Shafer2008, Vovk2005}. Miscoverage rate 0.1; 20% of each training fold held back for calibration. The base model is a 100-tree random forest. The score is the interval width, which scales with the spread of that forest's trees for the molecule.
% - Deep ensemble: 5 networks \citep{lakshminarayanan2017}. Spread of the member means = epistemic part; mean member variance = aleatoric part \citep{kendall2017}. The score combines both as one standard deviation.
% - Random forest uncertainty: the prose calls it "Bayesian random forest". What is computed is the spread across 200 trees plus half the out-of-fold absolute residual \citep{taddy2015}. The prose says out-of-bag; it is out-of-fold.
% - Quantile regression \citep{Koenker1978}: gradient boosting fitted at the 10th, 50th and 90th percentiles. The score is the width from the 10th to the 90th percentile.
% - Quantile regression forest \citep{Meinshausen2006}, missing from the prose: the score is how far the label lies outside its own predicted 10th-90th percentile range, 0 inside it. Unlike the width scores, it judges the label itself.
% - NGBoost predicts a mean and a standard deviation per molecule \citep{Duan2020}; the score is the standard deviation.
% - The last sentence of the prose is out of date in two ways. First, NGBoost is not selected: it is one score among the others. Second, the "promising" correlation it cites from Section~\ref{sec:uq} is pooled across noise levels; within one noise level it is near zero. Reword without "selected", and point to "Models with an explicit noise term track prediction error under noise", where those numbers are given once.
%
% Table: one row per noise score. Columns: family (nearest neighbours / out-of-fold models / models of their own uncertainty); what is computed, in one line; citation. It can replace the three itemised lists; keep the three formulas in the text.
%
% Decisions
% - Decision for you: whether the prose names the five scores that are computed but not yet described (distance-weighted neighbour residual, out-of-bag residual, combined score, heteroscedastic Gaussian process, quantile regression forest), or marks them as not reported. Naming them makes this list match the code; leaving them out keeps it to the scores whose results are reported.
% - This depends on the open decision in the parent subsection's notes: whether these estimator subsubsections stay in this chapter or move to the exploratory study's section. If they move, this whole list moves with them.
% ---- END NOTES ----

% ==== FROM: \subsubsection{Label smoothing and sample weighting} ====
% ---- NOTES: Label smoothing and sample weighting ----
% Paragraph 1: why move or down-weight labels instead of removing molecules
% - Every measured label carries some error, so removing a flagged molecule discards a label that is only partly wrong.
% - Smoothing moves a label part of the way toward a target value. Weighting keeps the label and lowers how much the molecule counts in training.
% - Both are noise mitigation: they act on the training labels.
% - The proportional and soft-threshold rules and all three weights need no yes/no call on which labels are noisy. The threshold rule does: it flags a share of molecules.
% - Define the two inputs once, before the lists. score_norm is the noise score from "Noise magnitude estimation", rescaled to 0-1 by its minimum and maximum over the fold's training molecules.
% - The target is the value a label is moved toward. It is the mean label of the molecule's nearest training neighbours, or a model's prediction for it (out-of-fold boosted trees, or a Gaussian-process posterior mean).
% - alpha is the strength, between 0 (no change) and 1 (label replaced by the target).
%
% Paragraph 2: the four smoothing rules
% - Threshold rule: the prose leaves out how far a flagged label moves. The top alpha share of molecules by score is flagged. Each flagged label moves alpha of the way to its target: $(1-\alpha)y + \alpha\,\text{target}$.
% - So in the threshold rule alpha sets both how many labels move and how far. Say so in one sentence.
% - Soft threshold: the strength shown is then used as in the proportional rule, $y + \text{strength}\cdot(\text{target} - y)$. Add this second line.
% - Confidence rule: uses the raw predicted standard deviation in label units, not score_norm. The target is the same model's predicted mean.
% - All four rules are for continuous labels. None of them defines what happens to a class label.
%
% Paragraph 3: the three weight shapes
% - The soft weight is clipped to 0.1-1.0 first and rescaled to mean 1 after, so a final weight can exceed 1.
% - Out of date: "to preserve effective sample size" is wrong. Mean 1 keeps the total weight equal to the number of training molecules.
% - The effective sample size, $(\sum w)^2/\sum w^2$, still falls as weights spread. Reword to "so the total weight equals the number of training molecules".
% - For the inverse weight, alpha is a ratio: the most suspect molecule gets $1/(1+\alpha)$ of the least suspect one's weight before rescaling.
% - Weights reach only the downstream model's fit. The hybrid's feature ranking never sees them.
%
% Paragraph 4: which of these rules the reported experiments use
% - Of the three experiments in this chapter, only the cleaning of real labels uses them: the proportional rule and the inverse weight. Its settings are listed in "Methods for learning with noisy labels"; point there, do not repeat them.
% - That experiment's move of only the flagged labels is not the threshold rule above. It replaces each flagged label with its neighbours' mean outright, and the flagged share (1\%, 5\% or 20\% of molecules) is set separately from any strength.
% - On a classification dataset, that experiment replaces the proportional move with a vote: each flagged label becomes the majority class of its nearest training neighbours.
% - The threshold, soft-threshold and confidence rules and the exponential and soft weights are not used in this chapter's experiments.
% - They were part of the grid in the exploratory study (Section~\ref{sec:lnl_motivation}), which swept the strength over 0.1, 0.3, 0.5 and 0.7 and used all four rules and all three weights. The confirmation report's methods list the same strengths, the same four rules and the same three weights \citep{Punt2026confirmation}. That key is not in refs.bib yet; the draft entry is in the chapter's opening notes.
% - Decision for you: keep the four rules and three weights in full as the exploratory study's methods, or cut the unused five to one sentence that points to that study. Kept, this subsubsection stays as long as it is and most of it serves another chapter. Cut, only the proportional rule and the inverse weight keep their formulas here.
% - Where this subsubsection sits (here or with the exploratory study) is the same choice recorded under "Methods for learning with noisy labels"; it is not a separate decision.
%
% Table: optional, one row per rule. Columns: rule; formula; what alpha controls; used in this chapter's experiments (yes/no). Replaces the two itemize lists if used.
% - If the table is used, define score_norm once in the sentence above it (Paragraph 1's definition), so the formulas in the table can use the symbol bare.
% ---- END NOTES ----

% ==== FROM: \subsubsection{Recovery rate} ====
% ---- NOTES: Recovery rate ----
% Paragraph 1: what the recovery rate measures
% - Recovery is the share of the score lost to label noise that an LNL method gives back.
% - The same downstream model is fitted three times on the same folds.
% - The three fits use clean, noisy and treated training labels.
% - Test labels are never altered in any of the three fits.
% - Recovery = (score after treatment - noisy score) / (clean score - noisy score), as a percentage.
% - 100% means clean performance is restored. 0% means no change.
% - A negative value means the treatment did worse than leaving the noise in.
% Paragraph 2: how it is computed here (the prose below describes an older single-model protocol; rewrite it)
% - The prose uses one random forest on ECFP4, scored by R^2.
% - The experiments in this chapter score every downstream model in the dataset's own metric (RMSE, Spearman or AUROC).
% - The formula works for any metric. Top and bottom are differences in the same metric, so its direction cancels.
% - For RMSE, where lower is better, both differences are negative. Flipping the sign of both gives the same ratio.
% - Each downstream model is divided by its own clean-to-noisy gap, not a gap shared across models.
% - The clean, noisy and treated scores are each averaged over folds before the ratio is taken.
% - No recovery is reported when the gap is under 2% of the clean score.
% - A near-zero gap turns small score differences into meaningless percentages.
% - No recovery is reported either when the noisy score is no worse than the clean score. There is then no loss to give back.
% - For labels clipped at a detection limit, the noisy score comes from training on the clipped labels. The formula is otherwise the same.
% - Replace the informal first-person opening ("We have a standard RF...") with a plain description of the three fits.
% - When rewriting the formula, use \[ ... \] for the display equation instead of $$ ... $$.
% - Decision for you: two per-model recovery tables do not follow the rules above. Either the code is changed so every table follows them, or the thesis names the two exceptions.
%   - The per-model table of the experiment that repairs labels after detection has no 2% rule. It refuses a recovery only when the gap is under 1e-10 or the noise did not hurt.
%   - The per-model table of the experiment comparing noise distributions has neither rule. It refuses a recovery only when the gap is under 1e-10, whatever its sign.
%   - If the code is changed, the bullets above hold for every table. If not, those two tables can show very large percentages where the gap is small, and one can show recovery where noise did not hurt.
% Paragraph 3: how detection is scored (the first sentence of the prose is out of date)
% - Decision for you: this paragraph is about detection, not recovery, but sits under "Recovery rate". Three options.
%   - Give it its own short paragraph heading here, e.g. \paragraph{Detection scores}, keeping it in Methods.
%   - Move your sentence and these bullets to the start of "Results: noise detection and estimation".
%   - Rename this subsubsection, e.g. "Recovery rate and detection scores".
%   - "Results: noise detection and estimation" currently points readers to "Recovery rate" for these definitions. That pointer changes with whichever option you pick.
% - The prose says precision, recall and F1. The experiments report three numbers that need no fixed cut-off.
% - Average precision: area under the precision-recall curve over all cut-offs. This is the headline number.
% - Precision in the top k: of the k molecules ranked most suspicious, the share whose label was changed.
% - k is the number of changed molecules, so at this k precision equals recall.
% - AUROC: the chance that a changed molecule scores above an unchanged one. 0.5 is a coin flip at any share of changed labels.
% - Each is reported next to the number of changed molecules and the number of training molecules.
% - Detection is not reported under a noise type that changes every label: no unchanged molecules exist to compare against.
% - Earlier exploratory cleaning-by-neighbours results are cited only in "Noise detection by neighbours" (Chapter 4). No table from that study is repeated here.
% ---- END NOTES ----

% ==== FROM: \section{Results: distances and nearest-neighbour prediction} ====
% ---- NOTES: Results: distances and nearest-neighbour prediction ----
% This section answers the first research question listed in this chapter's Introduction. Methods are defined in "Similarity in a hybrid" and "Feature importance weighting"; do not re-define them here.
%
% Paragraph 1: nearest-neighbour prediction in a hybrid against the single representations
% - This paragraph answers the research question directly.
% - Prediction: the mean activity of the k nearest training molecules, all counted equally, by Euclidean distance. k is 5 and 10 molecules.
% - Spaces compared: each of the dataset's single representations, the top three ordinary hybrids, and the top three similarity hybrids ("Similarity in a hybrid").
% - Each hybrid is scored twice: on its z-scored features, and with each feature scaled by its importance ("Feature importance weighting").
% - Budgets: 10, 25, 50, 100, 250 and 500 features.
% - Every space is scored in the dataset's own metric, on the held-out molecules of each fold. The neighbour score only decides which similarity hybrids enter.
% - Reference: for each held-out fold, the best single representation and the best hybrid are both chosen on the other folds. Report hybrid minus that representation.
% - The three datasets excluded in "Similarity in a hybrid" have no similarity hybrid, so no similarity-hybrid entry.
% - Report per-fold scores and their spread. No significance tests.
% - TODO(results): held-out score per dataset and per fold for the best single representation, the ordinary hybrid and the similarity hybrid, and each hybrid minus the representation.
% Table: one row per dataset. Columns: best single representation (its name and score), ordinary hybrid, similarity hybrid, each minus the single representation. Each entry is the mean over folds with the fold range.
%
% Paragraph 2: distance against activity difference
% - For sampled pairs of held-out molecules, plot the absolute difference in activity against their distance.
% - Use the same layout as the hERG Ki figure planned in "Flaws in similarity and other pairwise computations" (Chapter \ref{chap:flaws}), which plots activity difference against similarity for molecule pairs, so the two can be compared directly.
% - Two panels: the dataset's best single representation, and the similarity hybrid.
% - One number per panel: the Spearman correlation between pairwise distance and activity difference.
% - TODO(results): the figure, and the correlation per space for each dataset.
% Figure: two scatter panels side by side, activity difference (dataset units) against distance. Left: best single representation. Right: similarity hybrid. Caption gives the correlation for each panel.
%
% Paragraph 3: importance-weighted distance for k-nearest-neighbour prediction
% - Question: does scaling each feature by its importance give neighbours with closer activities than the unscaled features?
% - The comparison is defined in "Feature importance weighting" (Methods of this chapter, second paragraph of its notes). Point there; do not repeat the setup.
% - One setup fact the reader needs here: the space is all of the dataset's representations joined side by side, not a hybrid.
% - Its k-nearest-neighbour prediction uses Euclidean distance, where multiplying each feature by the square root of its importance matches "Feature importance weighting".
% - Its second prediction, from the 10 nearest training molecules by Manhattan distance, uses the same square-root scaling. Under Manhattan distance that multiplies each term by the square root of the importance, not by the importance itself as "Feature importance weighting" defines.
% - Decision for you: report the Manhattan comparison as well, or only the Euclidean one. Reporting it needs the convention decision in "Feature importance weighting" settled first; reporting only the Euclidean one needs no change to the code or the definition.
% - The score is R$^2$ on held-out molecules, not the dataset's own metric used in Paragraph 1. Say so in one sentence, unless the metric decision in "Feature importance weighting" changes it.
% - The comparison has no control that keeps the same weights but shuffles them across features. So weighted minus unweighted cannot say whether a gain comes from which features get the large weights, or only from the weights being unequal.
% - Decision for you: add that shuffled-weights control before the full runs, or report weighted against unweighted only. Added, it is one more scored space per fold and importance method, and a third column per k. Left out, the paragraph says the comparison cannot separate the two.
% - The importance-in-training comparison ("Results: why the hybrid helps") does not predict with nearest neighbours. It uses neighbours only to down-weight molecules whose label disagrees with theirs. It adds nothing to this paragraph.
% - TODO(results): unweighted and weighted R$^2$ per dataset, importance method, k and fold.
% Table: one row per dataset and importance method. Columns: unweighted R$^2$, weighted R$^2$ and their difference, at k = 5, 10 and 20 molecules.
%
% Paragraph 4: the similarity hybrid as an ordinary model input
% - In "Real labels, nothing injected: cleaning", downstream models are fitted on the similarity hybrid as well as on the ordinary hybrid and each model's top two single representations.
% - With no cleaning applied, this compares the two hybrids as inputs to the same downstream model, on the same folds.
% - The similarity hybrid's features were not chosen for a downstream model. State this as the reason the comparison is made, not as a result.
% - Scores are in each dataset's own metric, as in Chapter \ref{chap:method}.
% - TODO(results): held-out score with no cleaning, per dataset, downstream model and fold, for the similarity hybrid, the ordinary hybrid and each model's top two single representations.
%
% Whole section
% - Decision for you: whether this section reports further distance results beyond the four paragraphs above; each one added is a paragraph and a TODO(results) line, and leaving them out keeps the section to this chapter's first research question.
% - Candidates whose experiments are already written: distance from a held-out molecule to the training molecules (applicability domain), and prediction error on activity-cliff molecules.
% - Also: whether k-means clusters in the space hold molecules of similar activity, and retrieval of similar-activity molecules; Bayesian optimisation in a small hybrid.
% - Also: importance-weighted distances inside a Gaussian process and an SVM, and features built from each molecule's neighbours' activities.
% ---- END NOTES ----
% (agent-edited \todo, reverted to the 9763e83 wording)
\todo[inline, color=yellow!40]{Results TODO(prelim): methods are final, numbers will be replaced by the server runs.}

% ==== FROM: \section{Results: noise detection and estimation} ====
% ---- NOTES: Results: noise detection and estimation ----
% Scope: two of the three tasks set apart in Methods. Finding which training labels are wrong (detection), and sizing how wrong they are (estimation). Whether acting on either improves the held-out score is the next section; point there.
%
% Paragraph 1: how detection and estimation are scored (short; point back, do not re-define)
% - Each noise score ranks the training molecules of a fold from most to least suspect. The scores are defined in "Noise magnitude estimation".
% - The three detection numbers (average precision as headline, precision in the top k, AUROC) are defined in "Recovery rate". Point there.
% - Detection is graded against the molecules the injection singled out: widened under outlier noise (1\%, 5\% or 10\% of molecules) and grouped-wider noise, clipped under censoring.
% - Under outlier noise the other molecules still get ordinary noise. At 10\% the widened tenth carries half the injected squared error. Say so when reading its numbers.
% - Under Gaussian and grouped-shifted noise every training label moves, so no detection number exists. Only estimation applies there.
% - Chance level is the share of training molecules singled out, e.g. about 0.1 at 10\% outlier noise. Report it beside every average precision.
% - Estimation is graded against the injected noise size, in label units, or against a size measured from repeat measurements.
% - Feature spaces: each downstream model's top two single representations, the hybrid, and the hybrid ranked on the neighbour score (Similarity in a hybrid). The hybrids' search settings came from the search on clean labels (caveat in Methods; point there, do not repeat it).
% - Decision for you: keep this setup here, or move it into the chapter's Methods. Methods now describes three experiments (clipped labels, the twelve LNL methods, cleaning real labels) and none of them covers how detection and estimation are run. Moving it puts every setup in Methods and leaves this paragraph as one pointer. Keeping it leaves Methods without the setup for two of the chapter's results sections.
%
% Paragraph 2: finding injected noisy labels, and whether the space matters
% - Question: does a neighbour-based score find the corrupted molecules more often when neighbours are found in a hybrid than in a single representation?
% - Five ways to find the neighbours are compared: the single representation as it is; the same features weighted by importance (defined in "Feature importance weighting"); forest proximity, one minus the share of trees that put two molecules in the same leaf; the hybrid; and the hybrid ranked on the neighbour score.
% - The same comparison for the scores that fit a model instead of finding neighbours (out-of-fold residual, ensemble disagreement, the uncertainty models). These do not depend on a distance, so only the input features change.
% - Report per noise condition and level: average precision per score and space, with chance level, and the spread across folds and noise draws. No significance test.
% - TODO(results): average precision per score, space and noise condition, with chance level.
% - Preliminary numbers from the local runs, to quote marked preliminary until the server runs replace them:
%   - Local run, preliminary. Antibody thermostability (Tresanco), outlier noise on 10\% of labels, one-hot representation, 2 folds x 3 noise draws, no hybrid. Average precision against the moved molecules: plain Euclidean 0.1905; TreeSHAP-weighted Euclidean 0.2292; forest proximity 0.2251. Chance is about 0.09 (about 50 moved of 537 training molecules). These three were checked against the saved results.
%   - Local run, preliminary. hERG Ki, outlier noise on 10\% of labels, one fold: ten single representations reach average precision 0.4036 to 0.4326, and every way of combining the ten falls inside that range. Only the run at the lower noise level was saved (0.25 to 0.30 there), so this range is unverified.
%
% Paragraph 3: finding real suspect labels, nothing injected
% - Curator flags: ChEMBL curators mark records they judge suspect \citep{Zdrazil2023}. On the hERG IC50 records in ChEMBL release 37, 301 of 12,008 records carry a flag, on 215 of the 9,802 molecules.
% - 293 of the 301 are flagged "outside typical range" and 8 "potential transcription error". The first kind is assigned because the value is extreme. So a score that reads how far a label sits from its neighbours is partly graded on the value it was given. State this limit next to the number.
% - The 8 transcription-error records are too few to report on their own.
% - Grade at two levels: per record, and per molecule (a molecule counts once, flagged if any of its records is).
% - Inter-laboratory disagreement: BindingDB holds the same compound measured against the same protein by different laboratories \citep{Gilson2016}. The spread of their values is a real, per-molecule noise size nobody injected.
% - The flag marks outlier-type errors only: out-of-range values and typing mistakes. A score that finds them is not shown to find any other kind of noise. Say so.
% - Inter-laboratory disagreement is graded per molecule: the Spearman correlation between a score and the spread of that molecule's measurements.
% - Transfer: order the scores by how much of the noise damage their cleanup recovered under injected noise (Recovery rate). Order them again by their correlation with real disagreement. Report the Spearman correlation between the two orderings, only where at least 5 scores appear in both.
% - TODO(results): average precision against curator flags, per score and space; correlation with inter-laboratory spread; rank correlation of the two orderings.
%
% Paragraph 4: estimating how large the noise is
% - Three estimates that need no cleaning. All are in label units squared (e.g. pKi squared).
% - From repeat measurements: a molecule's label is the mean of its $k$ measurements, so its label error variance is the per-measurement variance divided by $k$.
% - Regress each molecule's squared out-of-fold residual on $1/k$. The slope is the per-measurement noise variance; the intercept is the model's own error variance. Show the residuals binned by $k$ (1, 2, 3, 4, 5 or more).
% - Beside it, the same variance measured directly: the pooled variance between repeat measurements of the same molecule. Only the hERG IC50 and BindingDB datasets keep repeats.
% - A bias shared by every measurement of a molecule (one value copied between papers) does not shrink with $k$. It lands in the intercept, so the slope counts only noise that differs between records.
% - From near-duplicate pairs: take the closest 1\% of training pairs by distance. Half their mean squared label difference, in the closest distance bin, estimates the noise variance, with no model. It still holds some real activity difference, so it is an upper bound.
% - From the gamma test \citep{Stefansson1997}: for each $k$ from 1 to 10, take the mean squared distance from each molecule to its $k$-th nearest neighbour, and the mean of half their squared label difference. Fit a straight line through the 10 points. Its intercept estimates the noise variance. It also holds real activity differences the features cannot resolve, so it is an upper bound.
% - Check on injected Gaussian noise: in a single representation the distances do not depend on the labels. So the rise in the estimate over level 0 should equal the injected variance.
% - Put each estimate beside the model's own held-out mean squared error. Noise mitigation spends model error to remove label error, so it can only pay where label error is a large share of the model's error.
% - Published size of Ki measurement error: already in Chapter 1 prose \citep{Kramer2012}. Point there and compare; do not restate it.
% - Forest proximity reads the labels when it groups molecules, so its near neighbours already agree. Its estimate is biased low by construction; report it beside the others, not instead of them.
% - TODO(results): the three estimates per dataset and space, beside the model's mean squared error.
%
% Paragraph 5: the gap between finding noise and correcting it
% - The point to state, if the results hold it up: finding which labels are noisy is possible for outlier noise, and the size of a systematic offset can be estimated. Correcting them to improve the held-out score is a different problem.
% - How precise a score must be: flag sets are built with their precision set by hand, from 0 to 1 in steps of 0.25, and each is acted on. This gives the precision a score would need before removing or down-weighting the flagged molecules improves the held-out score.
% - Compare it with the average precision the scores actually reach (Paragraph 2).
% - Link to "Results: learning with noisy labels, by noise type" for what acting on the flags does.
% - TODO(results): the precision needed before acting pays, per dataset, against the precision reached.
%
% Table: one row per noise score (grouped as in "Noise magnitude estimation"). Columns: average precision under outlier noise at 10\%, grouped-wider noise and censoring, each in the best single representation and in the hybrid; chance level in the header row.
% Figure: for one dataset, precision-recall curves of the best score in the best single representation, the importance-weighted single representation and the hybrid, with the chance line.
% ---- END NOTES ----

% ==== FROM: \section{Results: learning with noisy labels} ====
% ---- NOTES: Results: learning with noisy labels ----
% Paragraph 1: what this section reports, and in what order
% - Short lead-in, one paragraph. The results go experiment by experiment, then one closing subsection compares the noise types.
% - Three experiments, one subsection each: labels clipped at a detection limit; twelve LNL methods under Gaussian noise; cleaning real labels with nothing injected.
% - The first and third are noise mitigation: they change, drop or reweight training labels. The twelve-method survey covers both noise mitigation and noise-robust training.
% - The last subsection, By noise type, compares what each kind of noise allows. It holds the discussion of why Gaussian noise is the hardest case.
% - Methods, treatments and scores are defined in Methods for learning with noisy labels and Noise conditions. Point there; do not repeat them.
% - Outlier, grouped-wider and grouped-shifted noise have no result here. By noise type says so in one sentence; do not say it twice.
% - Four experiments on these noise types are written but have no result yet: correcting labels after detection under outlier noise; deleting against keeping flagged molecules; one offset per scaffold group; one offset per source document (hERG IC50, where each measurement names its paper).
% - Decision for you: whether any of these four goes into the first draft. If yes, each gets its own subsection here, after the real-label cleaning, and By noise type points to it instead of noting the gap. If no, this section stays at three experiments and By noise type keeps its one-sentence gap.
% - TODO(results): results of whichever of the four are kept.
%
% Paragraph 2: what every result in this section is compared against
% - Each downstream model's own top two single representations, the hybrid, and, where run, the importance-weighted hybrid or the similarity hybrid.
% - Scores are in each dataset's own metric: RMSE for hERG Ki and Caco-2, Spearman correlation for the antibody thermostability dataset.
% - Per-fold scores and their spread across folds are reported. No significance tests.
% - Caveat for the preliminary numbers in the subsections below (local runs, preliminary): they compare against one fixed single representation per dataset, one-hot on the antibody thermostability dataset and ECFP4 on hERG Ki and Caco-2.
% - One-hot is the weakest of the antibody dataset's ten representations for a random forest (0.3471 Spearman). Say this once here; the subsections point back.
% - TODO(results): every comparison against each model's top two single representations, replacing the preliminary numbers.
% ---- END NOTES ----

% ==== FROM: \subsection{Censoring: correcting labels clipped at a detection limit} ====
% ---- NOTES: Censoring: correcting labels clipped at a detection limit ----
% Paragraph 1: the question and why this noise type is studied by injection
% - Question: when training labels past an assay's limit are recorded at the limit, can they be corrected so the model predicts held-out molecules better?
% - The condition, the treatments and the recovery rate are defined in Noise conditions and Methods for learning with noisy labels. Point there; do not repeat the formulas.
% - Every treatment here is noise mitigation: it changes or drops the training labels.
% - Three models are comparators only: they change the model, not the labels. Each is compared with its parent model on the same clipped labels.
% - The three: a random forest with absolute-error splits, LightGBM with the Huber loss, and XGBoost with a censored likelihood.
% - The censored likelihood uses a normal error model. Each clipped label counts as a lower bound: the true value lies somewhere at or above the limit.
% - The treatments are not given the limit. They read it off the training labels as the largest recorded value.
% - Why injected, not real: the hERG Ki data keep only ChEMBL records with an exact "=" relation. Records reported as "<" or ">" a value never reach the experiment.
% - So the clean training labels are clipped on purpose. Held-out labels keep their real values.
% - Caco-2 is different: its loader already clips every label, training and held-out, at the dataset's 1st and 99th percentiles.
% - That clip is a data-loading choice, not an assay limit.
% - For Caco-2 at 25% clipped, the injected limit sits below the 99th percentile, so the injected clip covers the loader's upper clip. The 'never clipped' row still carries the loader's clip at both ends, about 1% of molecules each.
% - Whether the Caco-2 and LogD loaders keep this clip is one decision, asked once under Datasets and splitting in Chapter 5.
% - If kept, the held-out labels of both datasets stay censored at both ends.
% - If moved inside each training fold, the held-out labels stay as measured. The 'never clipped' row's training labels still carry a clip, with limits that differ from fold to fold.
% - If dropped, the labels of both datasets change in every experiment that uses them.
%
% Paragraph 2: preliminary result on hERG Ki at 10% clipped (local runs, preliminary)
% - Setup: ECFP4 only (2048 features), 3 folds, four downstream models (random forest, LightGBM, kNN, SVM).
% - The top 10% of training labels were set to 6.86 pKi: 88 to 100 of 941 or 942 training molecules per fold.
% - Correction shown: the constant lift, one amount added to every clipped label.
% - RMSE in pKi on held-out molecules, mean of 3 folds, lower is better. See the first table.
% - Recovery rate of the constant lift, per model: see the second table.
% - Deleting the clipped molecules scores worse than keeping them clipped, on all four models, in every fold and at every clipping level run.
% - Robust losses on the same clipped labels barely move from the clipped row of the first table.
% - Random forest with absolute-error splits: 0.6574 pKi. LightGBM with the Huber loss: 0.6480 pKi.
% - Both stay further from the never-clipped row than the constant lift does; see the first table.
% - Only ECFP4 was run, with a stand-alone preliminary script. The experiment in the final runs has not run yet; it uses each model's top two single representations and the hybrid.
% - The comparator with a censored likelihood (XGBoost, clipped labels as lower bounds) has no preliminary result.
%
% Paragraph 3: the other clipping levels and the per-molecule lift (local runs, preliminary)
% - Also run: hERG Ki at 25% (limit 6.20 pKi) and 40% (limit 5.81 pKi) clipped, and Caco-2 efflux ratio at 25% clipped.
% - Caco-2 scores are RMSE in log10 efflux ratio, 1438 or 1439 training molecules per fold.
% - The constant lift beat the clipped labels in 48 of 48 comparisons (4 models x 3 folds x 4 conditions).
% - Recovery shrinks as more labels are clipped on hERG Ki. See the second table.
% - A separate lift per clipped molecule scored higher RMSE than the single constant in 15 of 16 model-condition pairs.
% - The exception: LightGBM at 10%, 0.6245 against 0.6248 pKi.
% - Dealing the per-molecule lifts to the wrong clipped molecules (three random orders) scored lower RMSE than dealing them correctly, in 16 of 16 pairs.
% - Taken together: per-molecule lift sizes did no better than one shared size, and no better than the same sizes given to the wrong molecules.
% - Close by pointing to By noise type for why this noise type can be corrected and Gaussian noise cannot. Do not argue it here.
% - TODO(results): recovery rate per treatment, model, feature space and clipping level.
% - TODO(results): the three model comparators against their parent models, including XGBoost with the censored likelihood against XGBoost on the clipped and on the lifted labels.
%
% Table: hERG Ki, 10% of training labels clipped at 6.86 pKi, ECFP4, RMSE (pKi), mean of 3 folds. One row per version of the training labels.
% - Columns: random forest, LightGBM, kNN, SVM.
% - Never clipped: 0.6032, 0.6087, 0.6871, 0.6841.
% - Clipped: 0.6588, 0.6499, 0.7158, 0.7148.
% - Clipped, then constant lift: 0.6118, 0.6248, 0.6937, 0.6847.
% - Clipped molecules deleted: 0.7706, 0.7631, 0.8037, 0.8162.
% - TODO(results): replace with the final runs; add the top two single representations and the hybrid as feature spaces.
%
% Table: recovery rate of the constant lift, ECFP4, mean of 3 folds. One row per dataset and clipping level; columns random forest, LightGBM, kNN, SVM.
% - Recovery rate as defined in Recovery rate, with the clipped labels as the noisy fit. Do not restate the formula.
% - hERG Ki 10%: 85%, 61%, 77%, 98%.
% - hERG Ki 25%: 56%, 62%, 54%, 65%.
% - hERG Ki 40%: 45%, 48%, 51%, 53%.
% - Caco-2 25%: 65%, 71%, 72%, 68%.
% - Print the raw RMSE pairs beside each share, or in the appendix.
% ---- END NOTES ----

% ==== FROM: \subsection{Gaussian noise: learning-with-noisy-labels methods} ====
% ---- NOTES: Gaussian noise: learning-with-noisy-labels methods ----
% Paragraph 1: what was run (one short paragraph; the setup is in Methods)
% - Twelve LNL methods, each compared with one plain neural network trained with no noise handling. The methods and the network are defined in the methods table under Methods for learning with noisy labels. Do not describe them again.
% - Four of the twelve remove suspected molecules and retrain; these are noise mitigation. The other eight are noise-robust training.
% - Dataset in the preliminary runs: the antibody thermostability dataset, 672 antibodies, label is melting temperature.
% - Score: Spearman correlation between predicted and true melting temperature on test antibodies, higher is better, mean of 3 folds.
% - Gaussian noise added to training labels only, standard deviation 0.5 and 1.5 times the standard deviation of that fold's clean training labels. Levels are defined in Noise conditions.
% - Give the levels as these multiples, not in degrees Celsius. The value in Celsius differs by fold. The whole dataset's standard deviation is 10.12 C, so 0.5 and 1.5 times it would be 5.06 and 15.18 C; the per-fold values in the preliminary runs have not been checked.
% - In each fold the hybrid is built from that fold's noisy training labels, then the network trains on it.
% - Inputs in the preliminary runs: one-hot (5600 features), the hybrid, and the importance-weighted hybrid. The final runs use each model's top two single representations instead of one-hot.
%
% Paragraph 2: the plain network before any method (local runs, preliminary)
% - No added noise (not in the table, because the twelve methods were not run without noise): one-hot 0.4167, hybrid 0.5418, importance-weighted hybrid 0.5857 Spearman.
% - Noise at 0.5 times did not lower one-hot or the hybrid. Noise at 1.5 times did: the hybrid fell from 0.5418 with no added noise to 0.4859. The plain network's scores at both levels are the first row of the table.
% - So the 1.5 times level is the one that left damage for a method to repair.
%
% Paragraph 3: how many of the twelve methods beat the plain network on the same input (local runs, preliminary)
% - At 0.5 times: 0 of 12 on one-hot, 0 of 12 on the hybrid, 5 of 12 on the importance-weighted hybrid.
% - On the importance-weighted hybrid at 0.5 times the gains are small: the best method scores 0.5954 against 0.5833 for the plain network.
% - At 1.5 times: 8 of 12 on one-hot, 8 of 12 on the hybrid, 6 of 12 on the importance-weighted hybrid.
% - Removal methods at 1.5 times: small-loss selection beat the plain network on all three inputs. The other three removal methods beat it only on the importance-weighted hybrid (confident learning, 0.5316 against 0.5118), or not at all.
%
% Paragraph 4: input against method (local runs, preliminary)
% - The hybrid scored above one-hot at both noise levels, for the plain network and for every one of the twelve methods.
% - At 1.5 times, the best method on one-hot (small-loss selection, 0.3568) is below the plain network on the hybrid (0.4859).
% - At 1.5 times, the lowest of the twelve methods on the hybrid scores above the highest on one-hot, so the two ranges do not overlap. Point to the table's one-hot and hybrid columns at 1.5 times rather than restating the ranges.
% - State this as what was measured. The discussion of why noise-robust training is preferred is in Chapter~\ref{chap:discussion}; the discussion of why Gaussian noise is the hardest case is in By noise type. Point there.
%
% Paragraph 5: limits of these numbers
% - 3 folds and one noise draw per level.
% - In the preliminary runs the hybrid's features were searched with a random forest, while every method trains a neural network. The final runs search with the neural network and average three noise draws per fold.
% - The plain one-hot network ranges from 0.257 to 0.644 Spearman across folds with no added noise. That spread is larger than most gaps between methods.
% - One-hot is the weakest single representation on this dataset (see the lead-in to this section), so it is a weak comparator.
% - The final runs replace it with each model's top two single representations.
% - TODO(results): per noise level, the count of methods beating the plain network on each input, on every dataset in the final runs.
% - TODO(results): whether the input moves the score more than the method, with the final comparators.
%
% Table: one row per method (the plain network first, with 0.4534, 0.5991, 0.5833 at 0.5 times and 0.2729, 0.4859, 0.5118 at 1.5 times; then the twelve in the order of the methods table under Methods for learning with noisy labels); columns are one-hot, hybrid and importance-weighted hybrid at 0.5 times, then the same three at 1.5 times. Cells: Spearman, mean of 3 folds. A line above the table says what one row is.
% ---- END NOTES ----

% ==== FROM: \subsection{Real labels, nothing injected: cleaning} ====
% ---- NOTES: Real labels, nothing injected: cleaning ----
% Paragraph 1: what was run, in two sentences
% - The labels are used as they come; nothing is injected. Each cleanup is a noise-mitigation method: it removes, down-weights or moves training labels.
% - Design, suspicion scores, the seven cleanups and their settings are in Methods for learning with noisy labels. Point there; do not repeat them.
% - What a win can and cannot show is in Noise conditions. One sentence recalling it is enough: a better held-out score, not proof that wrong labels were found.
% - Datasets in the preliminary runs: hERG Ki and Caco-2 (RMSE, lower is better) and the antibody thermostability dataset (Spearman, higher is better).
% - Local runs, preliminary: random forest, 3 folds. Each number is the held-out score averaged over every setting of that cleanup.
% - The preliminary runs averaged over five suspicion scores. The final design uses two, plus the shuffled-score control. Every preliminary number below is replaced.
%
% Paragraph 2: the two small-molecule datasets
% - hERG Ki, RMSE in pKi units, local runs, preliminary. No cleaning: ECFP4 0.6032, hybrid 0.6029, similarity hybrid 0.6235.
% - On hERG Ki every cleanup, in every feature space, scored at or above its no-cleaning RMSE. Closest: down-weighting on ECFP4, 0.6040 against 0.6032.
% - Caco-2, RMSE in log efflux ratio, local runs, preliminary. No cleaning: ECFP4 0.3242, hybrid 0.3090, similarity hybrid 0.3226.
% - On Caco-2 down-weighting matched no cleaning: 0.3090 against 0.3090 on the hybrid, 0.3241 against 0.3242 on ECFP4. On the similarity hybrid it scored 0.3238 against 0.3226. Every other cleanup scored a higher RMSE in every feature space.
% - The Gaussian-process move was not run on Caco-2 in the preliminary runs.
% - Statement for the text: on both datasets, no cleanup improved on leaving the labels as they came.
%
% Paragraph 3: the antibody dataset, where the direction depends on the feature space
% - Spearman, local runs, preliminary. No cleaning: one-hot 0.4068, hybrid 0.6279, similarity hybrid 0.6087.
% - On one-hot every cleanup lowered the score. Worst: removal, 0.3498.
% - On the hybrid every cleanup raised the score. Largest: removal 0.6429, down-weighting 0.6407.
% - On the similarity hybrid the seven cleanups split: range clipping 0.6178 and two others above 0.6087; removal 0.6082 and three others below it.
% - Statement for the text: whether cleaning helped depended on the feature space, not only on the cleanup. Do not call it a gain until the shuffled-score control is in.
% - The shuffled-score control deals the same suspicion scores to random molecules. A gain that survives it came from which molecules were cleaned, not from how many.
% - The range clip reads no suspicion score, so the shuffled-score control cannot test it. A range-clip gain, such as 0.6178 Spearman on the similarity hybrid, has no control in this design.
% - TODO(results): the shuffled-score control per cleanup and feature space, on the antibody dataset first.
%
% Paragraph 4: what the final runs report
% - Headline: folds where the cleaned model beat the uncleaned one, out of folds run. Mean difference and spread across folds go beside it, never instead.
% - Feature spaces: each downstream model's top two single representations, the hybrid, and the similarity hybrid. The preliminary one-hot and ECFP4 columns are replaced by these.
% - Each downstream model is scored on the hybrid its own search picked, with that search's representations, allocation method and budget. No budget is fixed for this experiment, so do not quote one number of features for the hybrid.
% - QM9 is reported apart as the control. A gain there is not removal of measurement error.
% - The similarity hybrid's no-cleaning score as a downstream model input is discussed in Results: distances and nearest-neighbour prediction. Point there; do not discuss it here.
% - TODO(results): folds won per cleanup, suspicion score, feature space and dataset, including QM9.
%
% Table: one block per dataset. Rows: no cleaning, then the seven cleanups. Columns: each downstream model's top two single representations, the hybrid, the similarity hybrid. Cell: folds won out of folds run, with the mean difference in the dataset's metric. The line above the table says what one row is and what one cell counts.
% Table (Additional files, Chapter 6 group): the same, split by suspicion score and setting.
% ---- END NOTES ----

% ==== FROM: \subsection{By noise type} ====
% ---- NOTES: By noise type ----
% Paragraph 1: what this subsection does
% - Discussion, no new results. The numbers cited below are reported in earlier sections; point to them.
% - This is the one place that argues why Gaussian noise is the hardest case. Other sections point here.
% - Question for the subsection: what does each kind of noise allow you to find, and what does it allow you to fix?
% - Decision for you: this subsection's heading repeats the last words of its parent section, "Results: learning with noisy labels, by noise type", so the contents page lists "... by noise type" and then "By noise type" directly under it. Either rename this subsection (for example "What each kind of noise allows") or rename the section. Only the heading text and the pointers to it change.
% - The noise conditions used in this chapter are defined in "Noise conditions" in Methods. Recap each in one line at most. Use the names given there: Gaussian, censoring, outlier, grouped-wider and grouped-shifted.
%
% Paragraph 2: why Gaussian noise is the hardest to fix
% - Gaussian noise is the standard assumption in noise-robustness studies (Chapter~\ref{chap:background}, "Noise in molecular data"; point there). It is also the hardest to fix.
% - Every label carries an error of unknown size and unknown sign. No clean subset exists to learn from.
% - A detector that compares a molecule with its neighbours compares it with labels that are also noisy.
% - A large disagreement with the neighbours can be noise, or a real activity cliff.
% - Even a correctly flagged label has no known correction. Removing it loses a molecule. Replacing it with the neighbours' mean swaps its noise for theirs.
%
% Paragraph 3: finding the noise is not the same as removing it
% - The paradox: noise can be identified or estimated, yet a cleaned dataset can give a worse model than the noisy one it came from.
% - Evidence: in Chapter~\ref{chap:flaws}, "Noise detection by neighbours", the exploratory study ran 860 cleaning configurations on the QM9 HOMO-LUMO gap. Only 8 of the 860 beat training on the noisy labels at both noise levels.
% - Those two levels are Gaussian noise with standard deviation 0.3 eV and 0.6 eV, in the label's own units. They are not fractions of the label standard deviation, unlike the levels in this chapter.
% - This chapter's twelve-method survey does not show the same thing cleanly (local runs, preliminary). At 0.5 times the label standard deviation, none of the twelve beat the plain network on one-hot or on the hybrid. But noise at that level did not lower the plain network's score on either input, so there was no damage to repair.
% - At 1.5 times, 8 of the twelve beat the plain network on one-hot, and 8 of the twelve on the hybrid. Point to "Gaussian noise: learning-with-noisy-labels methods"; do not repeat its table.
% - So cite the exploratory study for the claim, and say the survey's preliminary numbers do not yet support it at either level.
% - TODO(results): whether the final survey shows any Gaussian level where the methods lose to the plain network after the noise has lowered its score.
% - A better similarity does not solve this. It can make neighbours closer in true activity. It cannot remove the noise on their labels.
%
% Paragraph 4: the other noise types, from easiest to fix to hardest
% - Censoring is the most fixable. The location (the limit) and the direction (up) are known; only the size is unknown.
% - That is why one constant lift recovered most of the clipping damage. Point to "Censoring: correcting labels clipped at a detection limit" for the numbers.
% - Outlier noise is easier to find than Gaussian noise. A few training molecules get a much wider error; the rest get ordinary noise, so the widened labels stick out against their neighbours.
% - The fix for an outlier is removal or down-weighting. Neither recovers the true value.
% - Grouped-wider noise gives some scaffold families a larger error of unknown sign. Every neighbour inside such a family is noisy too, so comparing a molecule with them says little.
% - Grouped-shifted noise moves each scaffold family by its own offset. Neighbours from the same family share the offset, so comparing with them cannot reveal it.
% - Knowing each record's source (laboratory, document) is what makes a shared offset estimable: fit one offset per source.
% - Outlier, grouped-wider and grouped-shifted noise have no LNL result in this chapter. Say so in one sentence; the parent section does not repeat it.
% - They do have detection and estimation results. Outlier and grouped-wider noise are used to grade detection; grouped-shifted noise only to grade estimation, because every training label moves under it. Point to "Results: noise detection and estimation".
% - Four experiments on outlier and grouped noise are written but not run. Whether they go into the first draft is decided in the notes of "Results: learning with noisy labels, by noise type". If they go in, point to their subsections here instead of the one-sentence gap.
%
% Paragraph 5: practical takeaways
% - First ask what kind of noise is likely, not only how much.
% - Data pooled from many laboratories or documents: expect an offset shared by each source, as in grouped-shifted noise. Keep the source of every record.
% - A single controlled experiment: expect censoring. Keep the "<" and ">" qualifiers instead of dropping those records.
% - Molecules from very different chemical domains: neighbours may not be close, so neighbour-based detection is least reliable.
% - Source and amount unknown: noise-robust training is safer than noise mitigation. The argument is made in Chapter~\ref{chap:discussion}, "Limitations"; point there.
% - TODO(results): per noise type, whether any LNL method, noise mitigation or noise-robust training, improved on the untreated model, on each feature space.
%
% Table: one row per noise type (Gaussian, outlier, grouped-wider, grouped-shifted, censoring).
% - Columns: what is wrong in the labels; how easy the noisy labels are to find; whether the direction of the error is known; what information makes a correction possible.
% - A line above the table says one row is one kind of label noise.
%
% Figure: five panels, one per noise type, in the table's order.
% - Each panel: x is the true label, y the observed label; the diagonal is no noise.
% - Shade the part that can be detected. Mark the part that can be corrected.
% - The figure shows why censoring can be fixed and Gaussian noise cannot.
% ---- END NOTES ----
% (agent-edited \todo, reverted to the 9763e83 wording)
\todo[inline, color=yellow!40]{Why Gaussian noise is hardest to fix; practical takeaways. Figure: five panels, one per noise type, as planned in the notes above.}

% ==== FROM: \section{Conclusions} ====
% ---- NOTES: Conclusions ----
% Short: 1-2 pages. One paragraph per research question, in the order of the Introduction, then noise type, limits and the bridge. No formulas; point back to the methods.
%
% Paragraph 1: distances and nearest-neighbour prediction (first research question)
% - Question: is predicting each molecule's activity from its five nearest neighbours more accurate in a hybrid than in the single representations it is built from?
% - Spaces to report: the ordinary hybrid and the similarity hybrid, each with plain and with importance-weighted distances.
% - Also say what the similarity hybrid scores as an ordinary downstream-model input, since its features were not chosen for a model.
% - TODO(results): per dataset, the held-out score of nearest-neighbour prediction in each space minus the best single representation, in the dataset's own metric.
% - TODO(results): per dataset, the similarity hybrid's held-out score as a downstream-model input.
% - If nearest-neighbour prediction improves when features are chosen for the property, then "similar" depends on the property, not on the two structures alone. Link back to Chapter \ref{chap:flaws}. Write this only if the results show it.
%
% Paragraph 2: noise detection and estimation (second research question)
% - Question: are noisy labels easier to find, or their size easier to estimate, in a hybrid?
% - State plainly that finding a noisy label and correcting it are separate problems. A method can do the first well and fail to improve the held-out score.
% - TODO(results): per dataset, detection and estimation in the hybrid against the top two single representations.
%
% Paragraph 3: learning with noisy labels (third research question)
% - Three experiments, one sentence each, in the order of the results sections.
% - Noise mitigation under censoring: TODO(results): whether correcting labels clipped at a detection limit beat leaving them, and how deleting them compared.
% - Learning with noisy labels under Gaussian noise: TODO(results): whether any of the twelve LNL methods beat the plain neural network at each noise level.
% - TODO(results): whether the choice of input moved the score more than the choice of method.
% - Noise mitigation on real labels, nothing injected: TODO(results): whether cleaning improved on leaving the labels, per dataset and feature space.
%
% Paragraph 4: what each noise type allows (two or three sentences; the argument is in "By noise type", do not repeat it)
% - The type of noise decides whether it can be corrected, not only found.
% - Censoring has a known location and direction. Gaussian noise has neither.
%
% Paragraph 5: limits specific to this chapter
% - The hybrid's features are rebuilt in each fold from that fold's training labels. So feature choice sees the noisy, clipped or uncleaned labels, as a real user's would.
% - In the clipped-label experiment the search settings (representations, budget, allocation method) come from the search on the unclipped labels. A real user would not have that head start.
% - The survey of twelve LNL methods re-searches the settings on the noisy labels, so this limit does not apply there.
% - The neighbour score is undefined on two of the 16 datasets: the breast cancer PAM50 subtype (4 classes) and trastuzumab-HER2 binding (a held-out fold can hold one group only). The reasons are in "Similarity in a hybrid"; point there. These two have no similarity hybrid. Their single representations and ordinary hybrids are still reported.
% - Decision for you: whether the two nucleotide datasets (H3K4me3, TATA promoter) are in the similarity results. The code searches the similarity hybrid on them like any other dataset. Keep them, and the similarity results cover 14 of the 16 datasets; drop them, and they cover 12 and this bullet must give the reason.
%
% Paragraph 6: bridge to Chapter \ref{chap:discussion}
% - Model the noise before choosing a method.
% - Which method helps depends on noise type, amount of noise, dataset and feature space.
% ---- END NOTES ----

% ==== FROM: \subsubsection{Noise magnitude estimation} (9763e83 comment, not carried into the chapter) ====
% TODO: [CR 863] carries the author's own '% TODO: get sources' marker. Left in place.

% ==== FROM: \subsubsection{Recovery rate} (9763e83 comments, not carried into the chapter) ====
% TODO: NOT carried over from [CR 985-1146], the confirmation report's hybrid results:
% TODO:   hybrid construction outcomes (tab:hybrid_both), smoothing-vs-weighting
% TODO:   (tab:smooth_vs_weight), the estimator table (tab:estimators), the per-representation
% TODO:   neighbour-based results, and the alpha sweep.
% TODO: The outline says the second-half results are pending a full re-run (Sec 5.6), and the
% TODO:   confirmation report's numbers came from the cache that was later invalidated. Recover from
% TODO:   Confirmation_Report/paper.tex lines 985-1146 if any of it survives the re-run.
% TODO: Also not carried: [CR 958-970], the QM9 hybrid-selection protocol (5,000 samples,
% TODO:   sigma in {0.3,0.6}, PDV+MHG-GNN budget=100 greedy allocation) -- superseded by the
% TODO:   ch.5 search; and [CR 1156], which is a cross-reference to the old chapter numbering.

% ==== Headings reverted to 9763e83 in pass 1 (agent renames, not private-name removals) ====
% \subsection{Methods for learning with noisy labels}  ->  \subsection{Noise mitigation methods}
% \section{Results: distances and nearest-neighbour prediction}  ->  \section{Results: similarity}
% \section{Results: learning with noisy labels}  ->  \section{Results: noise mitigation by noise type}

% ==== FROM: \paragraph{Feature importance weighting.} (integrity check after pass 1) ====
% Agent-added marker above her unchanged paragraph; the paragraph matches 9763e83, so the marker was false.
% TODO: review, needs manual review
